% Generated draft: requires mathematical and source review.
\documentclass[11pt]{article}
\usepackage[margin=25mm]{geometry}
\usepackage{fontspec}
\setmainfont{Latin Modern Roman}
\usepackage{amsmath,amssymb,mathtools,unicode-math}
\setmathfont{Latin Modern Math}
\usepackage{enumitem,needspace,xurl,hyperref}
\setlength{\parindent}{0pt}
\setlength{\parskip}{5pt}
\setlength{\emergencystretch}{3em}
\newcommand{\N}{\mathbb{N}}
\newcommand{\Z}{\mathbb{Z}}
\newcommand{\Q}{\mathbb{Q}}
\newcommand{\R}{\mathbb{R}}
\newcommand{\C}{\mathbb{C}}
\newcommand{\F}{\mathbb{F}}
\newcommand{\A}{\mathbb{A}}
\newcommand{\PP}{\mathbb{P}}
\newcommand{\E}{\mathbb{E}}
\newcommand{\eps}{\varepsilon}
\newcommand{\dd}{\,\mathrm{d}}
\begin{document}
\section*{UnsolvedMath ID: 20000319}
\textbf{Source problem code:} catalogue entry
\subsection{Definitions and mathematical statement}
Let \(\Delta=\{t\in\C:|t|<1\}\), and interpret a holomorphic family as a proper holomorphic submersion \(\pi:\mathcal X\to\Delta\) with compact connected K\"ahler fibers \(X_t\) of complex dimension \(n\). For every integer \(m\geq1\), define the plurigenus
\[
P_m(t)=\dim_{\C}H^0(X_t,K_{X_t}^{\otimes m}).
\]
Siu's question asks whether \(P_m(t)\) is independent of \(t\) for every \(m\). The original assumes K\"ahler fibers, not explicitly a K\"ahler total space. It points to the smooth projective invariance theorem, source reference [7], and explains the difficulty that an auxiliary ample line bundle used there need not be available in the K\"ahler case.

The attached aid reports the cases \(m=1\), fiber dimensions at most two, and projective or fiberwise Moishezon families as settled. It additionally reports a March 2026 Hacon--Li--Rao preprint establishing invariance for smooth compact K\"ahler threefold families, while describing fiber dimensions at least four as still open in general. This is its dated literature assessment, without publication details or new source verification here.

For fixed \(m\), write \(h_m^q(t)=\dim H^q(X_t,K_{X_t}^{\otimes m})\). The aid reports constancy of the Euler characteristic in a smooth proper family and the exact jump identity relative to generic fiber values:
\[
P_m(t)-P_m(\mathrm{gen})=
\sum_{q=1}^n(-1)^{q+1}\bigl(h_m^q(t)-h_m^q(\mathrm{gen})\bigr).
\]
Thus a special plurigenus jump requires that alternating higher-cohomology jump. If \(h_m^q(t_0)=0\) for all \(q>0\) on one fiber, the aid concludes local constancy of \(P_m\) near \(t_0\), not automatically constancy on the entire disk.

A line bundle is pseudoeffective in the compact K\"ahler setting when its first Chern class has a closed positive-current representative. The aid conditionally assumes deformation invariance of pseudoeffectivity of \(K_{X_t}\). It cites a proper holomorphic submersion with K\"ahler total space as a setting where this follows from results attributed to Fujiki--Levine and Ou. In such a family the non-pseudoeffective branch has
\[
P_m(t)=0\quad\text{for every }m\geq1,
\]
reducing its unresolved part to the pseudoeffective chamber. The additional total-space hypothesis and exact scope of the named theorems are not supplied in full, so the reduction is not silently extended to every fiberwise K\"ahler family.
\subsection{Short English statement}
Are all plurigenera constant in smooth proper families of compact K\"ahler manifolds? The attached aid reports known lower-dimensional cases and a March 2026 threefold preprint, then gives a higher-cohomology jump identity and a conditional reduction to pseudoeffective canonical bundles, leaving its stated higher-dimensional gap.
\subsection{Sources}
\begin{itemize}
\item Local UnsolvedMath registry, permanent ID \(20000319\), exact public code \texttt{catalogue entry}; canonical location \texttt{aim-algebraic-geometry-notes.json}, notes[318].
\item AIM workshop ``Numerical invariants of singularities and higher-dimensional algebraic varieties,'' attributed to Siu; supplied URL \url{https://aimath.org/WWN/singularvariety/problems_final.pdf}. The corpus prints Problem 9; its verification report corrects this to Problem 19 using the PDF and TeX source, and repairs K\"ahler typography. Neither source was consulted in this drafting step.
\item The original cites Siu's projective invariance theorem as [7], without the bibliography entry in this record. The original literature field is empty.
\item Attached research attempt 2, ``K\"ahler invariance of plurigenera: current boundary and a pseudoeffective-chamber reduction,'' mentioning a March 2026 Hacon--Li--Rao preprint and results of Fujiki--Levine and Ou without full bibliographic details.
\end{itemize}
\end{document}
